\chapter*{ઓપન લોજિક પ્રોજેક્ટ વિશે}
\addcontentsline{toc}{chapter}{ઓપન લોજિક પ્રોજેક્ટ વિશે}

\textit{ઓપન લોજિક ટેક્સ્ટ} રૂપલક્ષી અધિતર્કશાસ્ત્ર અને
રૂપલક્ષી પદ્ધતિઓનું મુક્ત સ્રોતવાળું સહયોગી પાઠ્યપુસ્તક છે.
તે મધ્યમ સ્તરથી શરૂ થાય છે, એટલે કે રૂપલક્ષી તર્કશાસ્ત્રનો
પ્રારંભિક અભ્યાસક્રમ પૂરો કર્યા પછીના સ્તરથી.
તે ગણિતની વિશેષ પૃષ્ઠભૂમિ ન ધરાવતા વાચકો માટે,
ખાસ કરીને તત્વજ્ઞાન અને સંગણકવિજ્ઞાનના વિદ્યાર્થીઓ માટે,
રચાયેલું હોવા છતાં તેની રજૂઆત કડક છે.

સમાવેલા કેટલાક વિષયોનું આવરણ હજી અધૂરું હોઈ શકે,
અને અનેક વિભાગોમાં હજી મોટા સુધારાની જરૂર છે.
ભવિષ્યમાં વધુ વિષયો આવરી લેવાનો અમારો ઇરાદો છે.
પાઠમાં પરિભાષાકોશ, વધુ વાંચનની સૂચિ, ઐતિહાસિક નોંધો,
ચિત્રો, વધુ સારી સમજણ, પરિણામોની તત્વજ્ઞાન,
સંગણકવિજ્ઞાન અને ગણિત માટેની પ્રાસંગિકતા સમજાવતા વિભાગો,
તથા વધુ કસરતો અને ઉદાહરણો ઉમેરવાની પણ યોજના છે.
ભૂલ મળે અથવા સૂચન હોય, તો
\href{https://github.com/OpenLogicProject/OpenLogic/wiki/Contributing}{પ્રોજેક્ટની ટીમને જણાવો}.

પ્રોજેક્ટ મુક્ત સ્રોતની ભાવનાથી ચાલે છે.
પાઠ મુક્તપણે ઉપલબ્ધ છે એટલું જ નહીં,
અમે તેનો LaTeX સ્રોત ક્રિએટિવ કૉમન્સ એટ્રિબ્યુશન પરવાના હેઠળ આપીએ છીએ.
યોગ્ય શ્રેય આપવાની શરતે કોઈ પણ વ્યક્તિને અમારું કાર્ય
ડાઉનલોડ કરવાનો, વાપરવાનો, બદલવાનો, ફરી ગોઠવવાનો,
અન્ય સ્વરૂપમાં ફેરવવાનો અને ફરી વિતરિત કરવાનો અધિકાર મળે છે.
વધુ માહિતી માટે
\href{http://openlogicproject.org/}{ઓપન લોજિક પ્રોજેક્ટનું જાળસ્થાન} જુઓ.

